% =====================================================================
% Section 7 -- ALTERNATE (v2): matched-vs-generic / motion-coverage narrative.
% A/B candidate against 07_interventions.tex. Main results table is
% tables/tab_realtargets.tex; the appearance gallery lives in the supplement.
% Swap by editing main.tex -> % =====================================================================
% Section 7 -- ALTERNATE (v2): matched-vs-generic / motion-coverage narrative.
% A/B candidate against 07_interventions.tex. Main results table is
% tables/tab_realtargets.tex; the appearance gallery lives in the supplement.
% Swap by editing main.tex -> % =====================================================================
% Section 7 -- ALTERNATE (v2): matched-vs-generic / motion-coverage narrative.
% A/B candidate against 07_interventions.tex. Main results table is
% tables/tab_realtargets.tex; the appearance gallery lives in the supplement.
% Swap by editing main.tex -> % =====================================================================
% Section 7 -- ALTERNATE (v2): matched-vs-generic / motion-coverage narrative.
% A/B candidate against 07_interventions.tex. Main results table is
% tables/tab_realtargets.tex; the appearance gallery lives in the supplement.
% Swap by editing main.tex -> \input{sections/07_interventions_v2}
% =====================================================================
\section{Interventions: Motion Coverage and Visual Realism}
\label{sec:interv}

Section~\ref{sec:law} found that the motion coverage of a source is a strong
predictor of its transfer to a target. We examine this relationship through
controlled experiments on two procedural generators with very different rendering.
One is Kubric~\cite{greff2021kubric}, a controllable generator that produces
photorealistic scenes from scanned assets and physically-based lighting. The other,
SDF-Fractals3D, is asset-free, built from procedural implicit surfaces, and shares
no objects, materials, or rendering with Kubric nor any benchmark datasets. We run two studies. The first is a
controlled grid that crosses five source-motion magnitudes with four appearance
variants over a fixed set of scenes, separating motion from appearance on KITTI-12 and KITTI-15. The
second compares, for each generator, a \emph{generic} target-agnostic configuration
against a \emph{motion-coverage-tuned} one across several real benchmarks and architectures.
The generic source is each generator's strongest single configuration. For Kubric, it
is the standard MOVi-F setup~\cite{greff2021kubric}. For SDF-Fractals3D, it is the best
configuration from a large generate-and-train grid search over the generator, found
independently of any coverage objective. The motion-coverage-tuned source instead
fits the generator's motion to the target using a search over motion coverage, training nothing. Across both studies, matching the source motion
to the target generally improves transfer, and the gains follow motion rather than
appearance.

\noindent\textbf{Motion magnitude drives transfer; appearance has little effect.}
This controlled grid uses Kubric, whose generator lets us re-render the same content
under different motion and appearance. We hold a fixed set of 5{,}000 scenes, each
with its own objects and layout, and vary only two factors across that set, the
source camera-motion magnitude across five settings and the visual appearance across
four variants spanning background, lighting, and object style. The scene content is
identical everywhere, so only these two axes change. Training a model on each of the
twenty resulting sources and evaluating on KITTI yields the transfer surface in
Fig.~\ref{fig:ladder:heatmaps}. Along the motion axis, transfer peaks where the source motion
best covers the target's and drops by $5$ to $11$ \pck{} as the magnitude moves away
from that match in either direction (Fig.~\ref{fig:ladder:heatmaps}), following the
coverage gap while off-target motion stays flat (Fig.~\ref{fig:ladder:coverage}).
Along the appearance axis, at identical motion, transfer is nearly flat. The most
photorealistic variant offers no consistent advantage over flat-shaded primitives,
even though its appearance distance to KITTI changes by roughly a factor of two
(Fig.~\ref{fig:ladder:heatmaps}). In this grid, transfer tracks coverage of the
target motion and varies little with appearance.

\begin{figure}[tbp]
\centering
\begin{subfigure}[t]{0.37\linewidth}
  \centering
  \includegraphics[height=3.0cm]{figures/results/ladder_panel_a.png}
  \caption{Directed motion distance ($\dbt$, $\dtb$) vs.\ source magnitude.}
  \label{fig:ladder:coverage}
\end{subfigure}%
\hfill%
\begin{subfigure}[t]{0.61\linewidth}
  \centering
  \includegraphics[height=3.0cm]{figures/results/ladder_panel_bc.png}
  \caption{Transfer to KITTI-2012 and KITTI-2015 (peak \pck{}).}
  \label{fig:ladder:heatmaps}
\end{subfigure}

\vspace{3pt}
\begin{subfigure}[b]{\linewidth}
  \centering
  \includegraphics[width=\linewidth]{figures/results/ladder_panel_d.png}
  \caption{Coverage gap ($\dbt$) resolved across the image frame, per source magnitude.}
  \label{fig:ladder:spatial}
\end{subfigure}
\caption{\textbf{Controlled motion and appearance grid evaluated on KITTI.}
Columns scale source motion magnitude to sweep coverage; rows vary appearance.
\emph{(a)} Coverage gap ($\dbt$, blue; lower is better) dips at the match and grows
on either side, while off-target motion ($\dtb$, red) stays flat.
\emph{(b)} \pck{} tracks coverage (columns) and is nearly constant across
appearance (rows). It peaks at the match and falls on either side, while higher
visual fidelity gives no consistent advantage. The rows
span about a factor of two in appearance distance to KITTI (\textbf{bold}).
\emph{(c)} The same coverage gap resolved across the KITTI frame, each region
colored relative to the best coverage it reaches across the sweep (green is that
best, red is worse). The match is best frame-wide and degrades on either side, most
in the high-magnitude lower half of the motion flow field. The per-variant appearance breakdown is shown in supplementary Fig.~\ref{S-fig:kubric_appearance}.}
\label{fig:ladder}
\end{figure}

\noindent\textbf{A geometry-only search finds an improved source.}
We search the generator's motion parameters with a Tree-structured Parzen Estimator
(TPE)~\cite{bergstra2011tpe}, using zero-shot motion coverage as the objective. Because the objective needs
only the candidate's and target's motion statistics, the search renders only lightweight
geometry passes, skipping the expensive texturing and lighting, and it trains no
model. The target's motion statistics come from a representative sample of the
benchmark. Only the selected configuration is later realized in full
(Fig.~\ref{fig:intervsplats}). We then train on each coverage-tuned source and compare
with the generic configuration. A construct-validity check that the coverage-tuned
parameters match each target's known motion regime is in the supplement
(Sect.~\ref{S-supp:recovery}).

\begin{figure}[tbp]
\centering
\includegraphics[width=\linewidth]{figures/results/F_pipeline_closedloop.png}
\caption{\textbf{Source design from motion alone.} A derivative-free search refines a
generic source with a single objective, motion coverage of the target, scored from
the motion field with no appearance rendered and no model trained. The coverage-tuned
source shows a coherent forward-motion field versus the generic source's
near-isotropic one, and training on it improves transfer (Table~\ref{tab:interv}).
Panels show directional motion fields; color $=$ direction.}
\label{fig:intervsplats}
\end{figure}

We evaluate the motion-coverage-tuned and generic sources on three real targets: KITTI and TSS
from the main study, and TAP-Vid-DAVIS~\cite{doersch2022tapvid}, a point-tracking
benchmark on real video held out from Sections~\ref{sec:setup}--\ref{sec:absolute}.
We read TAP-Vid-DAVIS at frame strides $s\in\{2,8,16\}$; a larger stride spaces the
two frames farther apart and increases the motion between them, so a single
benchmark spans small to large displacements.


We find that with SDF-Fractals3D the motion-coverage-tuned source improves on the generic one across nearly every
target, including the semantic benchmark TSS (Table~\ref{tab:interv}), and the gain is largest at the strict
KITTI threshold, where \pck{}@$1\%$ rises from at most $40$, the grid-searched
generic's best, to as high as $89$. With
Kubric the motion-coverage-tuned source improves the real-motion benchmarks, KITTI and
TAP-Vid-DAVIS, for most architectures. Its generic configuration is the standard
MOVi-F setup, a physics-based source whose broad motion already covers the semantic
target well, so on TSS the motion-coverage-tuned source has little room to improve and the two
stay close. This suggests that using motion coverage as an objective that predicts transfer can guide the design of a source datasets to improve transfer.

\begin{table}[tbp]
\centering
\caption{\textbf{Motion-coverage-tuned versus generic sources, in two generators.} Peak
\pck{} (\%) at discriminating thresholds: KITTI at $1\%$, TSS at $5\%$ and $3\%$,
and TAP-Vid-DAVIS at strides $s\in\{2,8,16\}$. The two generators are the
photorealistic Kubric generator (generic $=$ standard MOVi-F~\cite{greff2021kubric})
and our asset-free SDF-Fractals3D (generic $=$ best config from a large
generate-and-train grid search); \emph{MC-tuned} denotes the motion-coverage-tuned
configuration. \textbf{Bold} marks the better configuration
within each generator. Full thresholds and strides are in the supplement
(Tables~\ref{S-tab:interv_kitti_full} and~\ref{S-tab:interv_davis_full}).}
\label{tab:interv}
\setlength{\tabcolsep}{6pt}
\resizebox{\textwidth}{!}{\input{tables/tab_realtargets}}
\par\vspace{3pt}
\begin{minipage}{\textwidth}
\footnotesize $^{*}$\,SDF-generic on FlowFormer and RAFT did not converge and is
uncharacteristically poor, though the same source trains CATs++ and GLU-Net normally.
\end{minipage}
\end{table}

% =====================================================================
\section{Interventions: Motion Coverage and Visual Realism}
\label{sec:interv}

Section~\ref{sec:law} found that the motion coverage of a source is a strong
predictor of its transfer to a target. We examine this relationship through
controlled experiments on two procedural generators with very different rendering.
One is Kubric~\cite{greff2021kubric}, a controllable generator that produces
photorealistic scenes from scanned assets and physically-based lighting. The other,
SDF-Fractals3D, is asset-free, built from procedural implicit surfaces, and shares
no objects, materials, or rendering with Kubric nor any benchmark datasets. We run two studies. The first is a
controlled grid that crosses five source-motion magnitudes with four appearance
variants over a fixed set of scenes, separating motion from appearance on KITTI-12 and KITTI-15. The
second compares, for each generator, a \emph{generic} target-agnostic configuration
against a \emph{motion-coverage-tuned} one across several real benchmarks and architectures.
The generic source is each generator's strongest single configuration. For Kubric, it
is the standard MOVi-F setup~\cite{greff2021kubric}. For SDF-Fractals3D, it is the best
configuration from a large generate-and-train grid search over the generator, found
independently of any coverage objective. The motion-coverage-tuned source instead
fits the generator's motion to the target using a search over motion coverage, training nothing. Across both studies, matching the source motion
to the target generally improves transfer, and the gains follow motion rather than
appearance.

\noindent\textbf{Motion magnitude drives transfer; appearance has little effect.}
This controlled grid uses Kubric, whose generator lets us re-render the same content
under different motion and appearance. We hold a fixed set of 5{,}000 scenes, each
with its own objects and layout, and vary only two factors across that set, the
source camera-motion magnitude across five settings and the visual appearance across
four variants spanning background, lighting, and object style. The scene content is
identical everywhere, so only these two axes change. Training a model on each of the
twenty resulting sources and evaluating on KITTI yields the transfer surface in
Fig.~\ref{fig:ladder:heatmaps}. Along the motion axis, transfer peaks where the source motion
best covers the target's and drops by $5$ to $11$ \pck{} as the magnitude moves away
from that match in either direction (Fig.~\ref{fig:ladder:heatmaps}), following the
coverage gap while off-target motion stays flat (Fig.~\ref{fig:ladder:coverage}).
Along the appearance axis, at identical motion, transfer is nearly flat. The most
photorealistic variant offers no consistent advantage over flat-shaded primitives,
even though its appearance distance to KITTI changes by roughly a factor of two
(Fig.~\ref{fig:ladder:heatmaps}). In this grid, transfer tracks coverage of the
target motion and varies little with appearance.

\begin{figure}[tbp]
\centering
\begin{subfigure}[t]{0.37\linewidth}
  \centering
  \includegraphics[height=3.0cm]{figures/results/ladder_panel_a.png}
  \caption{Directed motion distance ($\dbt$, $\dtb$) vs.\ source magnitude.}
  \label{fig:ladder:coverage}
\end{subfigure}%
\hfill%
\begin{subfigure}[t]{0.61\linewidth}
  \centering
  \includegraphics[height=3.0cm]{figures/results/ladder_panel_bc.png}
  \caption{Transfer to KITTI-2012 and KITTI-2015 (peak \pck{}).}
  \label{fig:ladder:heatmaps}
\end{subfigure}

\vspace{3pt}
\begin{subfigure}[b]{\linewidth}
  \centering
  \includegraphics[width=\linewidth]{figures/results/ladder_panel_d.png}
  \caption{Coverage gap ($\dbt$) resolved across the image frame, per source magnitude.}
  \label{fig:ladder:spatial}
\end{subfigure}
\caption{\textbf{Controlled motion and appearance grid evaluated on KITTI.}
Columns scale source motion magnitude to sweep coverage; rows vary appearance.
\emph{(a)} Coverage gap ($\dbt$, blue; lower is better) dips at the match and grows
on either side, while off-target motion ($\dtb$, red) stays flat.
\emph{(b)} \pck{} tracks coverage (columns) and is nearly constant across
appearance (rows). It peaks at the match and falls on either side, while higher
visual fidelity gives no consistent advantage. The rows
span about a factor of two in appearance distance to KITTI (\textbf{bold}).
\emph{(c)} The same coverage gap resolved across the KITTI frame, each region
colored relative to the best coverage it reaches across the sweep (green is that
best, red is worse). The match is best frame-wide and degrades on either side, most
in the high-magnitude lower half of the motion flow field. The per-variant appearance breakdown is shown in supplementary Fig.~\ref{S-fig:kubric_appearance}.}
\label{fig:ladder}
\end{figure}

\noindent\textbf{A geometry-only search finds an improved source.}
We search the generator's motion parameters with a Tree-structured Parzen Estimator
(TPE)~\cite{bergstra2011tpe}, using zero-shot motion coverage as the objective. Because the objective needs
only the candidate's and target's motion statistics, the search renders only lightweight
geometry passes, skipping the expensive texturing and lighting, and it trains no
model. The target's motion statistics come from a representative sample of the
benchmark. Only the selected configuration is later realized in full
(Fig.~\ref{fig:intervsplats}). We then train on each coverage-tuned source and compare
with the generic configuration. A construct-validity check that the coverage-tuned
parameters match each target's known motion regime is in the supplement
(Sect.~\ref{S-supp:recovery}).

\begin{figure}[tbp]
\centering
\includegraphics[width=\linewidth]{figures/results/F_pipeline_closedloop.png}
\caption{\textbf{Source design from motion alone.} A derivative-free search refines a
generic source with a single objective, motion coverage of the target, scored from
the motion field with no appearance rendered and no model trained. The coverage-tuned
source shows a coherent forward-motion field versus the generic source's
near-isotropic one, and training on it improves transfer (Table~\ref{tab:interv}).
Panels show directional motion fields; color $=$ direction.}
\label{fig:intervsplats}
\end{figure}

We evaluate the motion-coverage-tuned and generic sources on three real targets: KITTI and TSS
from the main study, and TAP-Vid-DAVIS~\cite{doersch2022tapvid}, a point-tracking
benchmark on real video held out from Sections~\ref{sec:setup}--\ref{sec:absolute}.
We read TAP-Vid-DAVIS at frame strides $s\in\{2,8,16\}$; a larger stride spaces the
two frames farther apart and increases the motion between them, so a single
benchmark spans small to large displacements.


We find that with SDF-Fractals3D the motion-coverage-tuned source improves on the generic one across nearly every
target, including the semantic benchmark TSS (Table~\ref{tab:interv}), and the gain is largest at the strict
KITTI threshold, where \pck{}@$1\%$ rises from at most $40$, the grid-searched
generic's best, to as high as $89$. With
Kubric the motion-coverage-tuned source improves the real-motion benchmarks, KITTI and
TAP-Vid-DAVIS, for most architectures. Its generic configuration is the standard
MOVi-F setup, a physics-based source whose broad motion already covers the semantic
target well, so on TSS the motion-coverage-tuned source has little room to improve and the two
stay close. This suggests that using motion coverage as an objective that predicts transfer can guide the design of a source datasets to improve transfer.

\begin{table}[tbp]
\centering
\caption{\textbf{Motion-coverage-tuned versus generic sources, in two generators.} Peak
\pck{} (\%) at discriminating thresholds: KITTI at $1\%$, TSS at $5\%$ and $3\%$,
and TAP-Vid-DAVIS at strides $s\in\{2,8,16\}$. The two generators are the
photorealistic Kubric generator (generic $=$ standard MOVi-F~\cite{greff2021kubric})
and our asset-free SDF-Fractals3D (generic $=$ best config from a large
generate-and-train grid search); \emph{MC-tuned} denotes the motion-coverage-tuned
configuration. \textbf{Bold} marks the better configuration
within each generator. Full thresholds and strides are in the supplement
(Tables~\ref{S-tab:interv_kitti_full} and~\ref{S-tab:interv_davis_full}).}
\label{tab:interv}
\setlength{\tabcolsep}{6pt}
\resizebox{\textwidth}{!}{\begin{tabular}{ll rr rr rrr}
\toprule
 &  & \multicolumn{2}{c}{KITTI @1\%} & \multicolumn{2}{c}{TSS} & \multicolumn{3}{c}{TAP-Vid-DAVIS} \\
\cmidrule(lr){3-4}\cmidrule(lr){5-6}\cmidrule(lr){7-9}
Model & Source & K12 & K15 & @5\% & @3\% & $s{=}2$ & $s{=}8$ & $s{=}16$ \\
\midrule
\multirow{4}{*}{CATs++}
 & MOVi-F, generic & 45.9 & 46.2 & \textbf{61.1} & \textbf{44.1} & 98.5 & \textbf{89.7} & \textbf{84.0} \\
 & MOVi-F, MC-tuned & \textbf{46.1} & \textbf{46.8} & 60.9 & 43.5 & 98.5 & 89.5 & 83.2 \\
 & SDF, generic    & 40.2 & 44.5 & 67.0 & 49.8 & 98.1 & 89.4 & 83.2 \\
 & SDF, MC-tuned    & \textbf{89.1} & \textbf{76.3} & \textbf{67.1} & \textbf{51.0} & \textbf{98.4} & \textbf{90.6} & \textbf{84.1} \\
\addlinespace[3pt]
\multirow{4}{*}{GLU-Net}
 & MOVi-F, generic & 89.0 & 83.8 & 38.7 & 21.7 & 97.3 & 77.0 & 59.7 \\
 & MOVi-F, MC-tuned & \textbf{91.5} & \textbf{85.2} & \textbf{44.9} & \textbf{26.6} & \textbf{97.9} & \textbf{80.6} & \textbf{65.5} \\
 & SDF, generic    & 22.3 & 26.5 & 31.8 & 16.9 & 91.8 & 76.7 & \textbf{63.8} \\
 & SDF, MC-tuned    & \textbf{89.1} & \textbf{81.0} & \textbf{41.1} & \textbf{25.8} & \textbf{96.7} & \textbf{76.8} & 60.2 \\
\addlinespace[3pt]
\multirow{4}{*}{FlowFormer}
 & MOVi-F, generic & 58.4 & \textbf{56.1} & 30.1 & 16.9 & 95.2 & 65.7 & 48.8 \\
 & MOVi-F, MC-tuned & \textbf{59.8} & 53.2 & \textbf{32.9} & \textbf{18.5} & \textbf{95.6} & \textbf{68.3} & \textbf{51.0} \\
 & \textit{SDF, generic}$^{*}$ & \textit{6.2} & \textit{9.7} & \textit{32.7} & \textit{15.8} & \textit{93.3} & \textit{66.4} & \textit{50.1} \\
 & SDF, MC-tuned    & \textbf{62.7} & \textbf{56.0} & \textbf{37.3} & \textbf{22.7} & \textbf{96.0} & \textbf{70.4} & \textbf{52.0} \\
\addlinespace[3pt]
\multirow{4}{*}{RAFT}
 & MOVi-F, generic & \textbf{64.1} & 55.0 & \textbf{27.6} & \textbf{13.8} & 94.8 & 67.9 & 51.1 \\
 & MOVi-F, MC-tuned & 61.7 & 55.0 & 25.6 & 12.4 & \textbf{95.3} & \textbf{71.6} & \textbf{53.5} \\
 & \textit{SDF, generic}$^{*}$ & \textit{0.6} & \textit{0.6} & \textit{9.6} & \textit{3.5} & \textit{10.3} & \textit{12.6} & \textit{14.9} \\
 & SDF, MC-tuned    & \textbf{40.0} & \textbf{37.3} & \textbf{31.6} & \textbf{16.4} & \textbf{93.8} & \textbf{67.7} & \textbf{48.1} \\
\bottomrule
\end{tabular}
}
\par\vspace{3pt}
\begin{minipage}{\textwidth}
\footnotesize $^{*}$\,SDF-generic on FlowFormer and RAFT did not converge and is
uncharacteristically poor, though the same source trains CATs++ and GLU-Net normally.
\end{minipage}
\end{table}

% =====================================================================
\section{Interventions: Motion Coverage and Visual Realism}
\label{sec:interv}

Section~\ref{sec:law} found that the motion coverage of a source is a strong
predictor of its transfer to a target. We examine this relationship through
controlled experiments on two procedural generators with very different rendering.
One is Kubric~\cite{greff2021kubric}, a controllable generator that produces
photorealistic scenes from scanned assets and physically-based lighting. The other,
SDF-Fractals3D, is asset-free, built from procedural implicit surfaces, and shares
no objects, materials, or rendering with Kubric nor any benchmark datasets. We run two studies. The first is a
controlled grid that crosses five source-motion magnitudes with four appearance
variants over a fixed set of scenes, separating motion from appearance on KITTI-12 and KITTI-15. The
second compares, for each generator, a \emph{generic} target-agnostic configuration
against a \emph{motion-coverage-tuned} one across several real benchmarks and architectures.
The generic source is each generator's strongest single configuration. For Kubric, it
is the standard MOVi-F setup~\cite{greff2021kubric}. For SDF-Fractals3D, it is the best
configuration from a large generate-and-train grid search over the generator, found
independently of any coverage objective. The motion-coverage-tuned source instead
fits the generator's motion to the target using a search over motion coverage, training nothing. Across both studies, matching the source motion
to the target generally improves transfer, and the gains follow motion rather than
appearance.

\noindent\textbf{Motion magnitude drives transfer; appearance has little effect.}
This controlled grid uses Kubric, whose generator lets us re-render the same content
under different motion and appearance. We hold a fixed set of 5{,}000 scenes, each
with its own objects and layout, and vary only two factors across that set, the
source camera-motion magnitude across five settings and the visual appearance across
four variants spanning background, lighting, and object style. The scene content is
identical everywhere, so only these two axes change. Training a model on each of the
twenty resulting sources and evaluating on KITTI yields the transfer surface in
Fig.~\ref{fig:ladder:heatmaps}. Along the motion axis, transfer peaks where the source motion
best covers the target's and drops by $5$ to $11$ \pck{} as the magnitude moves away
from that match in either direction (Fig.~\ref{fig:ladder:heatmaps}), following the
coverage gap while off-target motion stays flat (Fig.~\ref{fig:ladder:coverage}).
Along the appearance axis, at identical motion, transfer is nearly flat. The most
photorealistic variant offers no consistent advantage over flat-shaded primitives,
even though its appearance distance to KITTI changes by roughly a factor of two
(Fig.~\ref{fig:ladder:heatmaps}). In this grid, transfer tracks coverage of the
target motion and varies little with appearance.

\begin{figure}[tbp]
\centering
\begin{subfigure}[t]{0.37\linewidth}
  \centering
  \includegraphics[height=3.0cm]{figures/results/ladder_panel_a.png}
  \caption{Directed motion distance ($\dbt$, $\dtb$) vs.\ source magnitude.}
  \label{fig:ladder:coverage}
\end{subfigure}%
\hfill%
\begin{subfigure}[t]{0.61\linewidth}
  \centering
  \includegraphics[height=3.0cm]{figures/results/ladder_panel_bc.png}
  \caption{Transfer to KITTI-2012 and KITTI-2015 (peak \pck{}).}
  \label{fig:ladder:heatmaps}
\end{subfigure}

\vspace{3pt}
\begin{subfigure}[b]{\linewidth}
  \centering
  \includegraphics[width=\linewidth]{figures/results/ladder_panel_d.png}
  \caption{Coverage gap ($\dbt$) resolved across the image frame, per source magnitude.}
  \label{fig:ladder:spatial}
\end{subfigure}
\caption{\textbf{Controlled motion and appearance grid evaluated on KITTI.}
Columns scale source motion magnitude to sweep coverage; rows vary appearance.
\emph{(a)} Coverage gap ($\dbt$, blue; lower is better) dips at the match and grows
on either side, while off-target motion ($\dtb$, red) stays flat.
\emph{(b)} \pck{} tracks coverage (columns) and is nearly constant across
appearance (rows). It peaks at the match and falls on either side, while higher
visual fidelity gives no consistent advantage. The rows
span about a factor of two in appearance distance to KITTI (\textbf{bold}).
\emph{(c)} The same coverage gap resolved across the KITTI frame, each region
colored relative to the best coverage it reaches across the sweep (green is that
best, red is worse). The match is best frame-wide and degrades on either side, most
in the high-magnitude lower half of the motion flow field. The per-variant appearance breakdown is shown in supplementary Fig.~\ref{S-fig:kubric_appearance}.}
\label{fig:ladder}
\end{figure}

\noindent\textbf{A geometry-only search finds an improved source.}
We search the generator's motion parameters with a Tree-structured Parzen Estimator
(TPE)~\cite{bergstra2011tpe}, using zero-shot motion coverage as the objective. Because the objective needs
only the candidate's and target's motion statistics, the search renders only lightweight
geometry passes, skipping the expensive texturing and lighting, and it trains no
model. The target's motion statistics come from a representative sample of the
benchmark. Only the selected configuration is later realized in full
(Fig.~\ref{fig:intervsplats}). We then train on each coverage-tuned source and compare
with the generic configuration. A construct-validity check that the coverage-tuned
parameters match each target's known motion regime is in the supplement
(Sect.~\ref{S-supp:recovery}).

\begin{figure}[tbp]
\centering
\includegraphics[width=\linewidth]{figures/results/F_pipeline_closedloop.png}
\caption{\textbf{Source design from motion alone.} A derivative-free search refines a
generic source with a single objective, motion coverage of the target, scored from
the motion field with no appearance rendered and no model trained. The coverage-tuned
source shows a coherent forward-motion field versus the generic source's
near-isotropic one, and training on it improves transfer (Table~\ref{tab:interv}).
Panels show directional motion fields; color $=$ direction.}
\label{fig:intervsplats}
\end{figure}

We evaluate the motion-coverage-tuned and generic sources on three real targets: KITTI and TSS
from the main study, and TAP-Vid-DAVIS~\cite{doersch2022tapvid}, a point-tracking
benchmark on real video held out from Sections~\ref{sec:setup}--\ref{sec:absolute}.
We read TAP-Vid-DAVIS at frame strides $s\in\{2,8,16\}$; a larger stride spaces the
two frames farther apart and increases the motion between them, so a single
benchmark spans small to large displacements.


We find that with SDF-Fractals3D the motion-coverage-tuned source improves on the generic one across nearly every
target, including the semantic benchmark TSS (Table~\ref{tab:interv}), and the gain is largest at the strict
KITTI threshold, where \pck{}@$1\%$ rises from at most $40$, the grid-searched
generic's best, to as high as $89$. With
Kubric the motion-coverage-tuned source improves the real-motion benchmarks, KITTI and
TAP-Vid-DAVIS, for most architectures. Its generic configuration is the standard
MOVi-F setup, a physics-based source whose broad motion already covers the semantic
target well, so on TSS the motion-coverage-tuned source has little room to improve and the two
stay close. This suggests that using motion coverage as an objective that predicts transfer can guide the design of a source datasets to improve transfer.

\begin{table}[tbp]
\centering
\caption{\textbf{Motion-coverage-tuned versus generic sources, in two generators.} Peak
\pck{} (\%) at discriminating thresholds: KITTI at $1\%$, TSS at $5\%$ and $3\%$,
and TAP-Vid-DAVIS at strides $s\in\{2,8,16\}$. The two generators are the
photorealistic Kubric generator (generic $=$ standard MOVi-F~\cite{greff2021kubric})
and our asset-free SDF-Fractals3D (generic $=$ best config from a large
generate-and-train grid search); \emph{MC-tuned} denotes the motion-coverage-tuned
configuration. \textbf{Bold} marks the better configuration
within each generator. Full thresholds and strides are in the supplement
(Tables~\ref{S-tab:interv_kitti_full} and~\ref{S-tab:interv_davis_full}).}
\label{tab:interv}
\setlength{\tabcolsep}{6pt}
\resizebox{\textwidth}{!}{\begin{tabular}{ll rr rr rrr}
\toprule
 &  & \multicolumn{2}{c}{KITTI @1\%} & \multicolumn{2}{c}{TSS} & \multicolumn{3}{c}{TAP-Vid-DAVIS} \\
\cmidrule(lr){3-4}\cmidrule(lr){5-6}\cmidrule(lr){7-9}
Model & Source & K12 & K15 & @5\% & @3\% & $s{=}2$ & $s{=}8$ & $s{=}16$ \\
\midrule
\multirow{4}{*}{CATs++}
 & MOVi-F, generic & 45.9 & 46.2 & \textbf{61.1} & \textbf{44.1} & 98.5 & \textbf{89.7} & \textbf{84.0} \\
 & MOVi-F, MC-tuned & \textbf{46.1} & \textbf{46.8} & 60.9 & 43.5 & 98.5 & 89.5 & 83.2 \\
 & SDF, generic    & 40.2 & 44.5 & 67.0 & 49.8 & 98.1 & 89.4 & 83.2 \\
 & SDF, MC-tuned    & \textbf{89.1} & \textbf{76.3} & \textbf{67.1} & \textbf{51.0} & \textbf{98.4} & \textbf{90.6} & \textbf{84.1} \\
\addlinespace[3pt]
\multirow{4}{*}{GLU-Net}
 & MOVi-F, generic & 89.0 & 83.8 & 38.7 & 21.7 & 97.3 & 77.0 & 59.7 \\
 & MOVi-F, MC-tuned & \textbf{91.5} & \textbf{85.2} & \textbf{44.9} & \textbf{26.6} & \textbf{97.9} & \textbf{80.6} & \textbf{65.5} \\
 & SDF, generic    & 22.3 & 26.5 & 31.8 & 16.9 & 91.8 & 76.7 & \textbf{63.8} \\
 & SDF, MC-tuned    & \textbf{89.1} & \textbf{81.0} & \textbf{41.1} & \textbf{25.8} & \textbf{96.7} & \textbf{76.8} & 60.2 \\
\addlinespace[3pt]
\multirow{4}{*}{FlowFormer}
 & MOVi-F, generic & 58.4 & \textbf{56.1} & 30.1 & 16.9 & 95.2 & 65.7 & 48.8 \\
 & MOVi-F, MC-tuned & \textbf{59.8} & 53.2 & \textbf{32.9} & \textbf{18.5} & \textbf{95.6} & \textbf{68.3} & \textbf{51.0} \\
 & \textit{SDF, generic}$^{*}$ & \textit{6.2} & \textit{9.7} & \textit{32.7} & \textit{15.8} & \textit{93.3} & \textit{66.4} & \textit{50.1} \\
 & SDF, MC-tuned    & \textbf{62.7} & \textbf{56.0} & \textbf{37.3} & \textbf{22.7} & \textbf{96.0} & \textbf{70.4} & \textbf{52.0} \\
\addlinespace[3pt]
\multirow{4}{*}{RAFT}
 & MOVi-F, generic & \textbf{64.1} & 55.0 & \textbf{27.6} & \textbf{13.8} & 94.8 & 67.9 & 51.1 \\
 & MOVi-F, MC-tuned & 61.7 & 55.0 & 25.6 & 12.4 & \textbf{95.3} & \textbf{71.6} & \textbf{53.5} \\
 & \textit{SDF, generic}$^{*}$ & \textit{0.6} & \textit{0.6} & \textit{9.6} & \textit{3.5} & \textit{10.3} & \textit{12.6} & \textit{14.9} \\
 & SDF, MC-tuned    & \textbf{40.0} & \textbf{37.3} & \textbf{31.6} & \textbf{16.4} & \textbf{93.8} & \textbf{67.7} & \textbf{48.1} \\
\bottomrule
\end{tabular}
}
\par\vspace{3pt}
\begin{minipage}{\textwidth}
\footnotesize $^{*}$\,SDF-generic on FlowFormer and RAFT did not converge and is
uncharacteristically poor, though the same source trains CATs++ and GLU-Net normally.
\end{minipage}
\end{table}

% =====================================================================
\section{Interventions: Motion Coverage and Visual Realism}
\label{sec:interv}

Section~\ref{sec:law} found that the motion coverage of a source is a strong
predictor of its transfer to a target. We examine this relationship through
controlled experiments on two procedural generators with very different rendering.
One is Kubric~\cite{greff2021kubric}, a controllable generator that produces
photorealistic scenes from scanned assets and physically-based lighting. The other,
SDF-Fractals3D, is asset-free, built from procedural implicit surfaces, and shares
no objects, materials, or rendering with Kubric nor any benchmark datasets. We run two studies. The first is a
controlled grid that crosses five source-motion magnitudes with four appearance
variants over a fixed set of scenes, separating motion from appearance on KITTI-12 and KITTI-15. The
second compares, for each generator, a \emph{generic} target-agnostic configuration
against a \emph{motion-coverage-tuned} one across several real benchmarks and architectures.
The generic source is each generator's strongest single configuration. For Kubric, it
is the standard MOVi-F setup~\cite{greff2021kubric}. For SDF-Fractals3D, it is the best
configuration from a large generate-and-train grid search over the generator, found
independently of any coverage objective. The motion-coverage-tuned source instead
fits the generator's motion to the target using a search over motion coverage, training nothing. Across both studies, matching the source motion
to the target generally improves transfer, and the gains follow motion rather than
appearance.

\noindent\textbf{Motion magnitude drives transfer; appearance has little effect.}
This controlled grid uses Kubric, whose generator lets us re-render the same content
under different motion and appearance. We hold a fixed set of 5{,}000 scenes, each
with its own objects and layout, and vary only two factors across that set, the
source camera-motion magnitude across five settings and the visual appearance across
four variants spanning background, lighting, and object style. The scene content is
identical everywhere, so only these two axes change. Training a model on each of the
twenty resulting sources and evaluating on KITTI yields the transfer surface in
Fig.~\ref{fig:ladder:heatmaps}. Along the motion axis, transfer peaks where the source motion
best covers the target's and drops by $5$ to $11$ \pck{} as the magnitude moves away
from that match in either direction (Fig.~\ref{fig:ladder:heatmaps}), following the
coverage gap while off-target motion stays flat (Fig.~\ref{fig:ladder:coverage}).
Along the appearance axis, at identical motion, transfer is nearly flat. The most
photorealistic variant offers no consistent advantage over flat-shaded primitives,
even though its appearance distance to KITTI changes by roughly a factor of two
(Fig.~\ref{fig:ladder:heatmaps}). In this grid, transfer tracks coverage of the
target motion and varies little with appearance.

\begin{figure}[tbp]
\centering
\begin{subfigure}[t]{0.37\linewidth}
  \centering
  \includegraphics[height=3.0cm]{figures/results/ladder_panel_a.png}
  \caption{Directed motion distance ($\dbt$, $\dtb$) vs.\ source magnitude.}
  \label{fig:ladder:coverage}
\end{subfigure}%
\hfill%
\begin{subfigure}[t]{0.61\linewidth}
  \centering
  \includegraphics[height=3.0cm]{figures/results/ladder_panel_bc.png}
  \caption{Transfer to KITTI-2012 and KITTI-2015 (peak \pck{}).}
  \label{fig:ladder:heatmaps}
\end{subfigure}

\vspace{3pt}
\begin{subfigure}[b]{\linewidth}
  \centering
  \includegraphics[width=\linewidth]{figures/results/ladder_panel_d.png}
  \caption{Coverage gap ($\dbt$) resolved across the image frame, per source magnitude.}
  \label{fig:ladder:spatial}
\end{subfigure}
\caption{\textbf{Controlled motion and appearance grid evaluated on KITTI.}
Columns scale source motion magnitude to sweep coverage; rows vary appearance.
\emph{(a)} Coverage gap ($\dbt$, blue; lower is better) dips at the match and grows
on either side, while off-target motion ($\dtb$, red) stays flat.
\emph{(b)} \pck{} tracks coverage (columns) and is nearly constant across
appearance (rows). It peaks at the match and falls on either side, while higher
visual fidelity gives no consistent advantage. The rows
span about a factor of two in appearance distance to KITTI (\textbf{bold}).
\emph{(c)} The same coverage gap resolved across the KITTI frame, each region
colored relative to the best coverage it reaches across the sweep (green is that
best, red is worse). The match is best frame-wide and degrades on either side, most
in the high-magnitude lower half of the motion flow field. The per-variant appearance breakdown is shown in supplementary Fig.~\ref{S-fig:kubric_appearance}.}
\label{fig:ladder}
\end{figure}

\noindent\textbf{A geometry-only search finds an improved source.}
We search the generator's motion parameters with a Tree-structured Parzen Estimator
(TPE)~\cite{bergstra2011tpe}, using zero-shot motion coverage as the objective. Because the objective needs
only the candidate's and target's motion statistics, the search renders only lightweight
geometry passes, skipping the expensive texturing and lighting, and it trains no
model. The target's motion statistics come from a representative sample of the
benchmark. Only the selected configuration is later realized in full
(Fig.~\ref{fig:intervsplats}). We then train on each coverage-tuned source and compare
with the generic configuration. A construct-validity check that the coverage-tuned
parameters match each target's known motion regime is in the supplement
(Sect.~\ref{S-supp:recovery}).

\begin{figure}[tbp]
\centering
\includegraphics[width=\linewidth]{figures/results/F_pipeline_closedloop.png}
\caption{\textbf{Source design from motion alone.} A derivative-free search refines a
generic source with a single objective, motion coverage of the target, scored from
the motion field with no appearance rendered and no model trained. The coverage-tuned
source shows a coherent forward-motion field versus the generic source's
near-isotropic one, and training on it improves transfer (Table~\ref{tab:interv}).
Panels show directional motion fields; color $=$ direction.}
\label{fig:intervsplats}
\end{figure}

We evaluate the motion-coverage-tuned and generic sources on three real targets: KITTI and TSS
from the main study, and TAP-Vid-DAVIS~\cite{doersch2022tapvid}, a point-tracking
benchmark on real video held out from Sections~\ref{sec:setup}--\ref{sec:absolute}.
We read TAP-Vid-DAVIS at frame strides $s\in\{2,8,16\}$; a larger stride spaces the
two frames farther apart and increases the motion between them, so a single
benchmark spans small to large displacements.


We find that with SDF-Fractals3D the motion-coverage-tuned source improves on the generic one across nearly every
target, including the semantic benchmark TSS (Table~\ref{tab:interv}), and the gain is largest at the strict
KITTI threshold, where \pck{}@$1\%$ rises from at most $40$, the grid-searched
generic's best, to as high as $89$. With
Kubric the motion-coverage-tuned source improves the real-motion benchmarks, KITTI and
TAP-Vid-DAVIS, for most architectures. Its generic configuration is the standard
MOVi-F setup, a physics-based source whose broad motion already covers the semantic
target well, so on TSS the motion-coverage-tuned source has little room to improve and the two
stay close. This suggests that using motion coverage as an objective that predicts transfer can guide the design of a source datasets to improve transfer.

\begin{table}[tbp]
\centering
\caption{\textbf{Motion-coverage-tuned versus generic sources, in two generators.} Peak
\pck{} (\%) at discriminating thresholds: KITTI at $1\%$, TSS at $5\%$ and $3\%$,
and TAP-Vid-DAVIS at strides $s\in\{2,8,16\}$. The two generators are the
photorealistic Kubric generator (generic $=$ standard MOVi-F~\cite{greff2021kubric})
and our asset-free SDF-Fractals3D (generic $=$ best config from a large
generate-and-train grid search); \emph{MC-tuned} denotes the motion-coverage-tuned
configuration. \textbf{Bold} marks the better configuration
within each generator. Full thresholds and strides are in the supplement
(Tables~\ref{S-tab:interv_kitti_full} and~\ref{S-tab:interv_davis_full}).}
\label{tab:interv}
\setlength{\tabcolsep}{6pt}
\resizebox{\textwidth}{!}{\begin{tabular}{ll rr rr rrr}
\toprule
 &  & \multicolumn{2}{c}{KITTI @1\%} & \multicolumn{2}{c}{TSS} & \multicolumn{3}{c}{TAP-Vid-DAVIS} \\
\cmidrule(lr){3-4}\cmidrule(lr){5-6}\cmidrule(lr){7-9}
Model & Source & K12 & K15 & @5\% & @3\% & $s{=}2$ & $s{=}8$ & $s{=}16$ \\
\midrule
\multirow{4}{*}{CATs++}
 & MOVi-F, generic & 45.9 & 46.2 & \textbf{61.1} & \textbf{44.1} & 98.5 & \textbf{89.7} & \textbf{84.0} \\
 & MOVi-F, MC-tuned & \textbf{46.1} & \textbf{46.8} & 60.9 & 43.5 & 98.5 & 89.5 & 83.2 \\
 & SDF, generic    & 40.2 & 44.5 & 67.0 & 49.8 & 98.1 & 89.4 & 83.2 \\
 & SDF, MC-tuned    & \textbf{89.1} & \textbf{76.3} & \textbf{67.1} & \textbf{51.0} & \textbf{98.4} & \textbf{90.6} & \textbf{84.1} \\
\addlinespace[3pt]
\multirow{4}{*}{GLU-Net}
 & MOVi-F, generic & 89.0 & 83.8 & 38.7 & 21.7 & 97.3 & 77.0 & 59.7 \\
 & MOVi-F, MC-tuned & \textbf{91.5} & \textbf{85.2} & \textbf{44.9} & \textbf{26.6} & \textbf{97.9} & \textbf{80.6} & \textbf{65.5} \\
 & SDF, generic    & 22.3 & 26.5 & 31.8 & 16.9 & 91.8 & 76.7 & \textbf{63.8} \\
 & SDF, MC-tuned    & \textbf{89.1} & \textbf{81.0} & \textbf{41.1} & \textbf{25.8} & \textbf{96.7} & \textbf{76.8} & 60.2 \\
\addlinespace[3pt]
\multirow{4}{*}{FlowFormer}
 & MOVi-F, generic & 58.4 & \textbf{56.1} & 30.1 & 16.9 & 95.2 & 65.7 & 48.8 \\
 & MOVi-F, MC-tuned & \textbf{59.8} & 53.2 & \textbf{32.9} & \textbf{18.5} & \textbf{95.6} & \textbf{68.3} & \textbf{51.0} \\
 & \textit{SDF, generic}$^{*}$ & \textit{6.2} & \textit{9.7} & \textit{32.7} & \textit{15.8} & \textit{93.3} & \textit{66.4} & \textit{50.1} \\
 & SDF, MC-tuned    & \textbf{62.7} & \textbf{56.0} & \textbf{37.3} & \textbf{22.7} & \textbf{96.0} & \textbf{70.4} & \textbf{52.0} \\
\addlinespace[3pt]
\multirow{4}{*}{RAFT}
 & MOVi-F, generic & \textbf{64.1} & 55.0 & \textbf{27.6} & \textbf{13.8} & 94.8 & 67.9 & 51.1 \\
 & MOVi-F, MC-tuned & 61.7 & 55.0 & 25.6 & 12.4 & \textbf{95.3} & \textbf{71.6} & \textbf{53.5} \\
 & \textit{SDF, generic}$^{*}$ & \textit{0.6} & \textit{0.6} & \textit{9.6} & \textit{3.5} & \textit{10.3} & \textit{12.6} & \textit{14.9} \\
 & SDF, MC-tuned    & \textbf{40.0} & \textbf{37.3} & \textbf{31.6} & \textbf{16.4} & \textbf{93.8} & \textbf{67.7} & \textbf{48.1} \\
\bottomrule
\end{tabular}
}
\par\vspace{3pt}
\begin{minipage}{\textwidth}
\footnotesize $^{*}$\,SDF-generic on FlowFormer and RAFT did not converge and is
uncharacteristically poor, though the same source trains CATs++ and GLU-Net normally.
\end{minipage}
\end{table}
